\documentclass[12pt,a4paper]{article}
\usepackage[utf8]{inputenc}
\usepackage[T1]{fontenc}
\usepackage[italian]{babel}
\usepackage{geometry}
\usepackage{amsmath,amssymb}
\usepackage{listings}
\usepackage{xcolor}
\usepackage{hyperref}
\usepackage{booktabs}
\usepackage{array}
\usepackage{tikz}
\usetikzlibrary{arrows.meta,positioning,shapes.geometric,fit,backgrounds,matrix}
\usepackage{float}
\usepackage{caption}
\usepackage{enumitem}
\usepackage{tcolorbox}
\tcbuselibrary{breakable,skins,theorems}
\usepackage{multicol}
\usepackage{mdframed}
\usepackage{fancyhdr}
\usepackage{titlesec}

\geometry{margin=2.2cm, top=2.5cm, bottom=2.5cm}

% ---- Header ----
\pagestyle{fancy}
\fancyhf{}
\fancyhead[L]{\small\textit{Intelligenza Artificiale -- Guida di Studio}}
\fancyhead[R]{\small\thepage}
\renewcommand{\headrulewidth}{0.4pt}

% ---- Colori ----
\definecolor{prologbg}{RGB}{248,250,255}
\definecolor{aspbg}{RGB}{250,255,250}
\definecolor{clipsbg}{RGB}{255,252,245}
\definecolor{pybg}{RGB}{245,248,245}
\definecolor{highlight}{RGB}{255,240,200}
\definecolor{keycolor}{RGB}{0,0,180}
\definecolor{commentcolor}{RGB}{60,120,60}
\definecolor{stringcolor}{RGB}{160,0,0}

% ---- Stile Prolog ----
\lstdefinelanguage{Prolog}{
  keywords={is,not,true,fail,halt,call,findall,member,memberchk,
            length,reverse,append,format,write,nl,assert,retract,assertz,
            retractall,forall,functor,arg,copy_term,nth0,nth1,
            atomic,is_list,maybe,random_permutation,statistics},
  keywordstyle=\color{keycolor}\bfseries,
  comment=[l]{\%},
  morecomment=[s]{/*}{*/},
  commentstyle=\color{commentcolor}\itshape,
  stringstyle=\color{stringcolor},
  morestring=[b]",
  sensitive=false,
  basicstyle=\ttfamily\footnotesize,
  backgroundcolor=\color{prologbg},
  frame=single,framerule=0.5pt,rulecolor=\color{blue!30},
  breaklines=true,showstringspaces=false,tabsize=4,
  numbers=left,numberstyle=\tiny\color{gray},numbersep=5pt,
}

% ---- Stile ASP ----
\lstdefinelanguage{ASP}{
  keywords={not},
  morekeywords=[2]{count,sum,min,max,minimize,maximize,show},
  keywordstyle=\color{keycolor}\bfseries,
  keywordstyle=[2]\color{purple!80!black}\bfseries,
  comment=[l]{\%},
  commentstyle=\color{commentcolor}\itshape,
  basicstyle=\ttfamily\footnotesize,
  backgroundcolor=\color{aspbg},
  frame=single,framerule=0.5pt,rulecolor=\color{green!40},
  breaklines=true,showstringspaces=false,tabsize=4,
  numbers=left,numberstyle=\tiny\color{gray},numbersep=5pt,
}

% ---- Stile CLIPS ----
\lstdefinelanguage{CLIPS}{
  keywords={deftemplate,defrule,defmodule,assert,retract,modify,focus,run,
            reset,slot,type,default,declare,salience,not,and,or,test},
  keywordstyle=\color{keycolor}\bfseries,
  comment=[l]{;;;},morecomment=[l]{;;},morecomment=[l]{;},
  commentstyle=\color{commentcolor}\itshape,
  basicstyle=\ttfamily\footnotesize,
  backgroundcolor=\color{clipsbg},
  frame=single,framerule=0.5pt,rulecolor=\color{orange!40},
  breaklines=true,showstringspaces=false,tabsize=4,
  numbers=left,numberstyle=\tiny\color{gray},numbersep=5pt,
}

\lstset{language=Prolog}

% ---- Box colorati ----
\newtcolorbox{defbox}[1]{colback=blue!6,colframe=blue!55!black,fonttitle=\bfseries\small,title=#1,breakable,left=4pt,right=4pt,top=3pt,bottom=3pt}
\newtcolorbox{exbox}[1]{colback=green!6,colframe=green!50!black,fonttitle=\bfseries\small,title=#1,breakable,left=4pt,right=4pt,top=3pt,bottom=3pt}
\newtcolorbox{warnbox}[1]{colback=orange!10,colframe=orange!75!black,fonttitle=\bfseries\small,title=#1,breakable,left=4pt,right=4pt,top=3pt,bottom=3pt}
\newtcolorbox{keybox}[1]{colback=yellow!15,colframe=yellow!60!black,fonttitle=\bfseries\small,title=#1,breakable,left=4pt,right=4pt,top=3pt,bottom=3pt}
\newtcolorbox{mathbox}[1]{colback=purple!6,colframe=purple!55!black,fonttitle=\bfseries\small,title=#1,breakable,left=4pt,right=4pt,top=3pt,bottom=3pt}

\hypersetup{colorlinks=true,linkcolor=blue!70!black,urlcolor=blue}

% ---- Titolo sezioni con colori ----
\titleformat{\section}{\large\bfseries\color{blue!70!black}}{{\thesection}}{0.8em}{}[\titlerule]
\titleformat{\subsection}{\normalsize\bfseries\color{blue!50!black}}{{\thesubsection}}{0.6em}{}
\titleformat{\subsubsection}{\small\bfseries\color{gray!80!black}}{{\thesubsubsection}}{0.5em}{}

% =====================================================================
\title{
  \vspace{-1em}
  \Large\textbf{\color{blue!70!black}Intelligenza Artificiale}\\[0.4em]
  \large\textbf{Guida Completa di Studio}\\[0.3em]
  \normalsize Prolog $\cdot$ ASP/Clingo $\cdot$ CLIPS $\cdot$ Algoritmi di Ricerca
}
\author{Alessandro Livero}
\date{\today}
% =====================================================================

\begin{document}
\maketitle
\tableofcontents
\newpage

% ======================================================================
% PARTE I - PROLOG
% ======================================================================
\part{Prolog -- Programmazione Logica}

% =====================================================================
\section{Fondamenti di Prolog}
% =====================================================================

\subsection{Cos'è Prolog}

Prolog (\textbf{Pro}gramming in \textbf{Log}ic) è un linguaggio di programmazione
\textbf{dichiarativo}: si descrive \emph{cosa} è vero, non \emph{come} calcolarlo.
Il motore di Prolog (l'interprete) si occupa di trovare le soluzioni tramite
\textbf{risoluzione SLD} (Selective Linear Definite clause resolution) e \textbf{backtracking}.

\subsection{I Tre Elementi Base}

\begin{center}
\renewcommand{\arraystretch}{1.4}
\begin{tabular}{llll}
\toprule
\textbf{Elemento} & \textbf{Sintassi} & \textbf{Significato} & \textbf{Esempio} \\
\midrule
Fatto & \texttt{pred(arg).} & Sempre vero & \texttt{gatto(tom).} \\
Regola & \texttt{h :- b1, b2.} & h è vero se b1 e b2 & \texttt{graffia(X) :- gatto(X).} \\
Query & \texttt{?- goal.} & Chiedi se è vero & \texttt{?- gatto(tom).} \\
\bottomrule
\end{tabular}
\end{center}

\begin{keybox}{Regola d'oro: lettere maiuscole = VARIABILI}
In Prolog tutto ciò che inizia con \textbf{lettera maiuscola} o \textbf{underscore \_}
è una \textbf{variabile}. Tutto il resto (lettere minuscole, numeri, simboli) è
una \textbf{costante} (atomo o numero).
\begin{lstlisting}
gatto(tom).      % tom e' una costante (atomo)
graffia(X) :- gatto(X).  % X e' una variabile
\end{lstlisting}
\end{keybox}

\subsection{Unificazione}

L'unificazione è il meccanismo fondamentale di Prolog: due termini si
\textbf{unificano} se esiste un'assegnazione di variabili che li rende identici.

\begin{exbox}{Esempi di unificazione}
\begin{lstlisting}
?- gatto(X) = gatto(tom).    % X = tom  (successo)
?- f(X, b) = f(a, Y).       % X = a, Y = b  (successo)
?- f(X) = g(X).             % fallisce (f != g)
?- X = f(X).                % X = f(f(f(...))) (occur check)
\end{lstlisting}
\end{exbox}

\subsection{Backtracking}

Quando una query fallisce o si vogliono più soluzioni, Prolog torna indietro
(\textbf{backtrack}) e prova la prossima clausola disponibile.

\begin{lstlisting}
uccello(titty).
pinguino(tux).
uccello(X) :- pinguino(X).   % un pinguino e' un uccello

vola(X) :- uccello(X), \+ pinguino(X).
\end{lstlisting}

\textbf{Traccia di} \texttt{?- vola(titty)}:
\begin{enumerate}[noitemsep]
  \item Cerca clausola con testa \texttt{vola(?)}: trova \texttt{vola(X) :- uccello(X), \textbackslash+ pinguino(X)}. Unifica X = titty.
  \item Prova \texttt{uccello(titty)}: successo (è un fatto).
  \item Prova \texttt{\textbackslash+ pinguino(titty)}: prova \texttt{pinguino(titty)} $\to$ fallisce $\to$ \textbackslash+ succede. \textbf{Risultato: vero.}
\end{enumerate}

\textbf{Traccia di} \texttt{?- vola(tux)}:
\begin{enumerate}[noitemsep]
  \item Unifica X = tux. Prova \texttt{uccello(tux)}: non è un fatto diretto.
  \item Prova regola: \texttt{uccello(X) :- pinguino(X)}, X=tux. \texttt{pinguino(tux)} $\to$ successo. Quindi \texttt{uccello(tux)} succede.
  \item Prova \texttt{\textbackslash+ pinguino(tux)}: \texttt{pinguino(tux)} $\to$ successo $\to$ \textbackslash+ \textbf{fallisce}. \textbf{Risultato: falso.}
\end{enumerate}

\subsection{Negazione per Fallimento (\texttt{\textbackslash+})}

\begin{defbox}{Negazione per Fallimento (NAF)}
\texttt{\textbackslash+ Goal} è vero se \texttt{Goal} \textbf{fallisce} (non è dimostrabile).
Implementa la \textbf{Closed World Assumption (CWA)}: ciò che non è derivabile
si assume falso.

\textbf{Non è} la negazione logica classica: \texttt{\textbackslash+ p(X)} non significa
``esiste una prova che p(X) è falso'', ma ``non riesco a provare p(X)''.
\end{defbox}

\subsection{Non-Monotonic Reasoning (NMR) -- file \texttt{nmr.pl}}

\begin{lstlisting}
% Il paradosso del pinguino
uccello(X) :- pinguino(X).
vola(X)    :- uccello(X), \+ pinguino(X).
uccello(titty).
pinguino(tux).
pinguino(titty).   % titty e' anche un pinguino!
\end{lstlisting}

\textbf{Il ragionamento non-monotonico}: normalmente gli uccelli volano.
Ma aggiungere il fatto \texttt{pinguino(titty)} \textbf{ritira} la conclusione
\texttt{vola(titty)}, anche se era stata derivata prima. Il sistema è
\emph{non-monotono}: aggiungere informazioni può eliminare conclusioni.

Questo è diverso dalla logica classica (monotona) dove aggiungere assiomi non
toglie mai teoremi.

% =====================================================================
\section{Liste in Prolog}
% =====================================================================

\subsection{Rappresentazione}

Una lista è o la lista vuota \texttt{[]} oppure una coppia \texttt{[Testa|Coda]}
dove \texttt{Testa} è il primo elemento e \texttt{Coda} è il resto (altra lista).

\begin{lstlisting}
[1, 2, 3]    % lista di 3 elementi
[H|T]        % H = 1, T = [2, 3]
[A, B | R]   % A = 1, B = 2, R = [3]
[]           % lista vuota
\end{lstlisting}

\subsection{Pattern: Ricorsione su Liste}

Il pattern universale per lavorare su liste in Prolog:
\begin{lstlisting}
predicato([]).               % caso base: lista vuota
predicato([Head|Tail]) :-    % passo ricorsivo
    ... Head ...             % fai qualcosa con la testa
    predicato(Tail).         % ricorsione sulla coda
\end{lstlisting}

\subsection{Predicati su Liste -- file \texttt{lista.pl}}

\subsubsection{\texttt{tuttiPositivi/1} -- Controllo Uniforme}
\begin{lstlisting}
tuttiPositivi([]).                         % lista vuota: ok
tuttiPositivi([Head|Tail]) :-
    atomic(Head), Head > 0,                % elemento > 0
    tuttiPositivi(Tail).
tuttiPositivi([Head|Tail]) :-
    is_list(Head),                         % se e' una sottolista
    tuttiPositivi(Head),                   % controlla la sottolista
    tuttiPositivi(Tail).
\end{lstlisting}

\textbf{Gestisce liste annidate}: \texttt{tuttiPositivi([1, [2, 3], 4])} $\to$ vero.
Il predicato \texttt{atomic/1} controlla che non sia una lista; \texttt{is\_list/1}
controlla che sia una lista.

\subsubsection{\texttt{ContaPositivi/2} -- Conteggio}
\begin{lstlisting}
ContaPositivi([], 0).
ContaPositivi([Head|Tail], Risultato) :-
    atomic(Head), Head > 0,
    ContaPositivi(Tail, Tot),
    Risultato is Tot + 1.
ContaPositivi([Head|Tail], Risultato) :-
    atomic(Head), Head =< 0,              % elemento <= 0: non conta
    ContaPositivi(Tail, Risultato).
ContaPositivi([Head|Tail], Risultato) :-
    is_list(Head),                         % sottolista: conta ricorsivamente
    ContaPositivi(Head, Tot1),
    ContaPositivi(Tail, Tot2),
    Risultato is Tot1 + Tot2.
\end{lstlisting}

\begin{warnbox}{Attenzione: \texttt{is} per l'aritmetica}
In Prolog \texttt{X = 3 + 2} \textbf{non calcola}: unifica X con il termine
\texttt{+(3,2)}. Per valutare, usare \texttt{X is 3 + 2} (X = 5).
\end{warnbox}

\subsection{Inversione di Lista -- file \texttt{invertilista.pl}}

\begin{lstlisting}
% Versione naive: O(n^2)
% inverti_naive([1,2,3], X) => X = [3,2,1]

% Versione efficiente con accumulatore: O(n)
inverti(Lista, LInv) :- inv(Lista, [], LInv).

inv([], Acc, Acc).               % lista esaurita: l'accumulatore E' il risultato
inv([H|T], Acc, LInv) :-
    inv(T, [H|Acc], LInv).      % metti H in testa all'accumulatore (inverti)
\end{lstlisting}

\textbf{Traccia di} \texttt{inv([1,2,3], [], Inv)}:
\begin{itemize}[noitemsep]
  \item[] \texttt{inv([1,2,3], [])} $\to$ \texttt{inv([2,3], [1])} $\to$ \texttt{inv([3], [2,1])} $\to$ \texttt{inv([], [3,2,1])} $\to$ \texttt{Inv = [3,2,1]}
\end{itemize}

\textbf{Il pattern dell'accumulatore} è fondamentale: invece di ricostruire la
lista durante il ritorno dalla ricorsione (lento), si costruisce il risultato
\emph{in avanti} passandolo come parametro aggiuntivo.

\subsection{Operazioni su Insiemi -- file \texttt{listeinsiemi.pl}}

\begin{lstlisting}
% Intersezione: elementi di L1 che stanno anche in L2
intersez([], _, []).
intersez(_, [], []).
intersez([H|T], L2, [H|R]) :- member(H, L2), !, intersez(T, L2, R).
intersez([_|T], L2, R)     :- intersez(T, L2, R).

% Unione: tutti gli elementi di L1 e L2 senza duplicati
unione([], L2, L2) :- !.
unione(L1, [], L1) :- !.
unione([H|T], L2, R)      :- member(H, L2), !, unione(T, L2, R).
unione([H|T], L2, [H|R])  :- unione(T, L2, R).

% Differenza: elementi di L1 non presenti in L2
diff([], _, []) :- !.
diff([H|T], L2, [H|R]) :- \+ member(H, L2), !, diff(T, L2, R).
diff([_|T], L2, R)     :- diff(T, L2, R).
\end{lstlisting}

% =====================================================================
\section{Il Cut (\texttt{!})}
% =====================================================================

\subsection{Cos'è il Cut}

Il cut (\texttt{!}) è un predicato speciale che, quando viene incontrato,
\textbf{congela} le scelte già fatte: elimina i punti di backtracking creati
dalla clausola corrente e dalle sue sorelle.

\begin{defbox}{Effetto del Cut}
Quando \texttt{!} viene eseguito nella clausola \texttt{h :- b1, !, b2}:
\begin{enumerate}[noitemsep]
  \item Tutte le clausole alternative per \texttt{h} vengono \textbf{scartate}.
  \item I punti di backtracking creati da \texttt{b1} vengono \textbf{eliminati}.
  \item Se \texttt{b2} fallisce, la query per \texttt{h} fallisce definitivamente
        (non prova le clausole successive).
\end{enumerate}
\end{defbox}

\subsection{Cut Verde vs Cut Rosso}

\begin{center}
\begin{tabular}{lll}
\toprule
\textbf{Tipo} & \textbf{Effetto sulla semantica} & \textbf{Uso} \\
\midrule
\textbf{Cut verde} & Non cambia il significato logico & Solo efficienza \\
\textbf{Cut rosso} & Cambia il significato (rimuove soluzioni) & Esprimere priorità \\
\bottomrule
\end{tabular}
\end{center}

\subsection{Esempio: \texttt{contaPositivi} con Cut -- file \texttt{listeconcut.pl}}

\begin{lstlisting}
contaPositivi([], 0).
contaPositivi([Head|Tail], Risultato) :-
    Head > 0, !,                     % se Head > 0, CUT: evita di provare le clausole sotto
    contaPositivi(Tail, Tot),
    Risultato is Tot + 1.
contaPositivi([_|Tail], Tot) :-      % questa clausola vale solo se Head <= 0
    contaPositivi(Tail, Tot).
\end{lstlisting}

\textbf{Senza cut}: Prolog proverebbe tutte e 3 le clausole per ogni elemento.
\textbf{Con cut}: se \texttt{Head > 0} riesce, si esegue solo la seconda clausola
(quella con il cut), e la terza non viene mai provata. Se \texttt{Head > 0} fallisce,
si salta alla terza clausola. Questo è un \textbf{cut rosso}: modifica il
comportamento (la terza clausola non si attiva mai quando Head > 0).

\begin{exbox}{Altro uso del cut: \texttt{prova/1}}
\begin{lstlisting}
prova(X) :- member(X, [1,2,3,4,5]), !, member(X, [4,5,10,50]).
prova(152).
\end{lstlisting}
\texttt{?- prova(X).} trova il primo X in [1,2,3,4,5] (= 1), poi il cut impedisce
di tornare sulla \texttt{member}. Allora prova \texttt{member(1, [4,5,10,50])} $\to$
fallisce. Il cut impedisce anche di provare la seconda clausola \texttt{prova(152)}.
Risultato: \textbf{fallisce}. Senza il cut otterrebbe X = 4 o X = 5.
\end{exbox}

\subsection{Esempio: \texttt{sottolista} con Cut}
\begin{lstlisting}
sottolista([Head|_], Head) :- is_list(Head), !.  % primo elemento lista: fermati
sottolista([_|Tail], R)   :- sottolista(Tail, R). % altrimenti continua
\end{lstlisting}
Il cut garantisce che quando si trova il primo elemento lista, non si cerchino
altre soluzioni nella stessa posizione.

% =====================================================================
\section{Ricorsione e Alberi Genealogici -- file \texttt{antenati.pl}}
% =====================================================================

\subsection{Predicati Genealogici}

\begin{lstlisting}
genitore(edoardo, gianni).  % ... altri fatti ...
genitore(gianni, margherita).
genitore(margherita, johnElkann).

nonno(X, Y) :- genitore(X, Z), genitore(Z, Y).
\end{lstlisting}

\textbf{Traccia di} \texttt{?- nonno(edoardo, margherita)}:
\begin{enumerate}[noitemsep]
  \item Unifica X=edoardo, Y=margherita.
  \item Cerca Z: \texttt{genitore(edoardo, Z)} $\to$ Z=gianni (primo risultato).
  \item Controlla: \texttt{genitore(gianni, margherita)} $\to$ successo!
\end{enumerate}

\subsection{Antenato: Ricorsione su Relazioni}

\begin{lstlisting}
% Attenzione: questa definizione ha un bug!
antenato(X, Y) :- antenato(X, Z), genitore(Z, Y).
\end{lstlisting}

\begin{warnbox}{Bug: ricorsione a sinistra infinita}
La definizione sopra va in loop infinito perché la prima chiamata a
\texttt{antenato} genera subito un'altra chiamata a \texttt{antenato}
senza caso base. La versione corretta dovrebbe essere:
\begin{lstlisting}
antenato(X, Y) :- genitore(X, Y).           % caso base
antenato(X, Y) :- genitore(X, Z), antenato(Z, Y).  % ricorsione
\end{lstlisting}
La ricorsione \textbf{a destra} (il predicato ricorsivo appare dopo una
chiamata a un fatto) termina correttamente.
\end{warnbox}

% =====================================================================
\section{Predicati Higher-Order: \texttt{findall}, \texttt{call}}
% =====================================================================

\subsection{\texttt{findall/3}}

\begin{defbox}{\texttt{findall(+Template, +Goal, -Lista)}}
Raccoglie in \texttt{Lista} tutti i valori di \texttt{Template} per cui
\texttt{Goal} è soddisfatto. Se \texttt{Goal} non ha soluzioni, \texttt{Lista = []}.

\textbf{Non dà mai errore}: se non trova niente, restituisce lista vuota.
\end{defbox}

\begin{lstlisting}
haTipo(pichu, elettro).
haTipo(emboar, fuoco).
haTipo(volcanion, fuoco).
haTipo(volcanion, acqua).

% ?- findall(P, haTipo(P, fuoco), Lista).
% Lista = [emboar, volcanion]

% ?- findall(P-T, haTipo(P, T), Coppie).
% Coppie = [pichu-elettro, emboar-fuoco, volcanion-fuoco, volcanion-acqua]
\end{lstlisting}

\subsection{\texttt{call/N} e Meta-predicati}

\texttt{call(Goal)} esegue \texttt{Goal} come se fosse scritto direttamente.
\texttt{call(Goal, Arg1)} aggiunge \texttt{Arg1} come argomento a \texttt{Goal}.

\begin{lstlisting}
:- meta_predicate a_star(1, 3, 2, +, -, -, -).
% significa: il primo arg e' un predicato che call chiamera' con 1 extra arg

% Uso: a_star(maze_goal, maze_successors, maze_heuristic, Start, ...)
% internamente chiama: call(maze_goal, State)
%                      call(maze_successors, State, NS, Cost)
%                      call(maze_heuristic, State, H)
\end{lstlisting}

% =====================================================================
\section{A* e IDA* in Prolog -- file \texttt{astar\_idastar.pl}}
% =====================================================================

\subsection{Il Problema della Ricerca nello Spazio degli Stati}

Dato uno \textbf{stato iniziale}, un insieme di \textbf{operatori} (mosse),
e una \textbf{condizione di goal}, trovare la sequenza di operatori che porta
dallo stato iniziale a uno stato goal con \textbf{costo minimo}.

\subsection{A* (A-Star)}

\subsubsection{Idea Chiave}
A* esplora sempre il nodo con il valore $f = g + h$ minimo:
\begin{itemize}[noitemsep]
  \item $g(n)$: costo reale dallo stato iniziale a $n$
  \item $h(n)$: stima euristica (ammissibile: $h(n) \le h^*(n)$) da $n$ al goal
  \item $f(n) = g(n) + h(n)$: stima del costo totale attraverso $n$
\end{itemize}

\subsubsection{Implementazione Prolog}

\begin{lstlisting}
% Struttura nodo nella frontiera: f(ValoreF, G, Stato, PercorsoInverso)
% PercorsoInverso e' la lista degli stati percorsi, in ordine inverso (efficiente)

astar(Domain, Start, Path, Cost, Expanded) :-
    heuristic(Domain, Start, H),
    % frontiera iniziale: un solo nodo con f=h, g=0
    astar_loop(Domain, [f(H, 0, Start, [Start])], [], Path, Cost, 0, Expanded).

% Caso base: il nodo in testa e' il goal -> SOLUZIONE TROVATA
astar_loop(Domain, [f(_, G, State, RevPath)|_], _, Path, G, Exp, Exp) :-
    goal(Domain, State),
    reverse(RevPath, Path).    % inverte il percorso costruito al contrario

% Caso ricorsivo: espansione
astar_loop(Domain, [f(_, G, State, RevPath)|Open], Closed, Path, Cost, Exp0, Exp) :-
    \+ goal(Domain, State),
    (member(State, Closed) ->           % gia' espanso: scarta (lazy deletion)
        astar_loop(Domain, Open, Closed, Path, Cost, Exp0, Exp)
    ;
        Exp1 is Exp0 + 1,
        findall(f(F1, G1, S1, [S1|RevPath]),
            (move(Domain, State, S1, SC),   % genera successori
             \+ member(S1, Closed),
             G1 is G + SC,
             heuristic(Domain, S1, H1),
             F1 is G1 + H1),
            Successors),
        insert_all(Successors, Open, Open1),  % inserimento ordinato per F
        astar_loop(Domain, Open1, [State|Closed], Path, Cost, Exp1, Exp)
    ).
\end{lstlisting}

\subsubsection{Perché il Goal si Controlla all'Estrazione?}

Si controlla il goal quando il nodo è \textbf{estratto} (in testa alla lista), non
quando viene inserito. Questo garantisce l'ottimalità: quando un nodo goal è in testa
alla frontiera (ordinata per $f$), il suo $g$ è il minimo possibile.

\subsection{IDA* (Iterative Deepening A*)}

\subsubsection{Idea Chiave}
Esegue DFS iterative con soglie $t$ crescenti su $f$:
\begin{enumerate}[noitemsep]
  \item Prima iterazione: $t_0 = h(\text{start})$
  \item DFS: espande nodi con $f \le t$; se $f > t$ ritaglia (cutoff)
  \item Se la DFS non trova il goal: $t_{\text{new}} = \min(f$ dei nodi ritagliati$)$
  \item Ripete finché trova il goal o nessun nodo è ritagliabile
\end{enumerate}

\textbf{Vantaggio}: memoria $O(b \cdot d)$ invece di $O(b^d)$ come A*.
\textbf{Svantaggio}: riesplora i nodi ad ogni iterazione.

\begin{lstlisting}
idastar(Domain, Start, Path, Cost, Iters) :-
    heuristic(Domain, Start, H0),
    idastar_loop(Domain, Start, H0, Path, Cost, 0, Iters).

idastar_loop(Domain, Start, Bound, Path, Cost, Iters0, Iters) :-
    idastar_search(Domain, Start, 0, Bound, [Start], ResultPath, ResultCost, NextBound),
    (ResultPath \= failure ->
        Path = ResultPath, Cost = ResultCost, Iters is Iters0 + 1
    ;
        Iters1 is Iters0 + 1,
        idastar_loop(Domain, Start, NextBound, Path, Cost, Iters1, Iters)
    ).

% Caso base: goal trovato
idastar_search(Domain, State, G, _, RevPath, Path, G, _) :-
    goal(Domain, State), reverse(RevPath, Path), !.

% Caso ricorsivo: cutoff o espansione
idastar_search(Domain, State, G, Bound, RevPath, Path, Cost, NextBound) :-
    heuristic(Domain, State, H), F is G + H,
    (F > Bound ->                          % CUTOFF: f supera la soglia
        Path = failure, Cost = 0, NextBound = F
    ;
        findall(s(S1, SC),
            (move(Domain, State, S1, SC),
             \+ member(S1, RevPath)),      % evita cicli (path checking)
            Succs),
        search_successors(Domain, Succs, G, Bound, RevPath, Path, Cost, 1.0e308, NextBound)
    ).
\end{lstlisting}

\subsection{Interfaccia Generica}

\begin{lstlisting}
% Collegamento domain -> predicati specifici
move(maze, S, S1, C)   :- maze_move(S, S1, C).
move(puzzle, S, S1, C) :- puzzle_move(S, S1, C).

heuristic(maze, S, H)   :- maze_heuristic(S, H).
heuristic(puzzle, S, H) :- puzzle_heuristic(S, H).

goal(maze, S)   :- maze_goal(S).
goal(puzzle, S) :- puzzle_goal_check(S).
\end{lstlisting}

\textbf{Aggiungere un nuovo dominio} richiede solo definire tre predicati:
\texttt{move(nuovo, ...)}, \texttt{heuristic(nuovo, ...)}, \texttt{goal(nuovo, ...)}.

\subsection{Dominio 1: Labirinto (5$\times$5)}

\begin{lstlisting}
% Stato: pos(Riga, Colonna). Muri: wall(R, C). Uscite: maze_exit(pos(R,C)).

maze_move(pos(R, C), pos(R1, C), 1) :-    % SU
    R1 is R - 1, maze_size(MaxR, _),
    R1 >= 1, R1 =< MaxR, \+ wall(R1, C).
% ... analoghe per GIU', SINISTRA, DESTRA

maze_heuristic(pos(R, C), H) :-            % Manhattan verso uscita piu' vicina
    findall(D, (maze_exit(pos(Re, Ce)), D is abs(R-Re)+abs(C-Ce)), Ds),
    min_list(Ds, H).
\end{lstlisting}

\textbf{Ammissibilità della Manhattan}: nei movimenti ortogonali con costo 1,
la Manhattan è il numero minimo di passi senza ostacoli. I muri possono solo
aumentare il costo reale, quindi $h \le h^*$ sempre.

\subsection{Dominio 2: Puzzle dell'8}

\textbf{Stato}: lista di 9 elementi \texttt{[0,1,2,3,4,5,6,7,8]}, lo 0 è la casella vuota.

\begin{lstlisting}
% Mapping indice -> (riga, colonna):
%   Riga = Indice div 3,   Colonna = Indice mod 3

puzzle_move(State, NewState, 1) :-
    blank_pos(State, Blank),
    (   Blank >= 3, Swap is Blank - 3   % SU
    ;   Blank =< 5, Swap is Blank + 3   % GIU'
    ;   Col is Blank mod 3, Col > 0, Swap is Blank - 1  % SINISTRA
    ;   Col is Blank mod 3, Col < 2, Swap is Blank + 1  % DESTRA
    ),
    swap_elements(State, Blank, Swap, NewState).

% Euristica: somma delle distanze Manhattan di tutte le tessere
puzzle_heuristic(State, H) :-
    puzzle_goal(Goal),
    findall(D,
        (nth0(Idx, State, Tile), Tile \= 0,
         nth0(GoalIdx, Goal, Tile),
         R1 is Idx // 3, C1 is Idx mod 3,
         R2 is GoalIdx // 3, C2 is GoalIdx mod 3,
         D is abs(R1-R2) + abs(C1-C2)),
        Ds),
    sumlist(Ds, H).
\end{lstlisting}

\begin{keybox}{Puzzle dell'8: condizione di risolvibilità}
Non tutte le permutazioni di [0..8] sono risolvibili! Una permutazione è
risolvibile \textbf{se e solo se il numero di inversioni è pari}.
Un'inversione è una coppia $(A, B)$ con $A$ prima di $B$, $A > B$, e nessuno
dei due è 0. Esattamente metà delle 362880 permutazioni è risolvibile.
\end{keybox}

\subsection{Confronto A* vs IDA*}

\begin{center}
\renewcommand{\arraystretch}{1.4}
\begin{tabular}{lll}
\toprule
\textbf{Proprietà} & \textbf{A*} & \textbf{IDA*} \\
\midrule
Memoria & $O(b^d)$ -- esponenziale & $O(b \cdot d)$ -- lineare \\
Nodi riesplorati & No & Sì (ogni iterazione) \\
Struttura dati & Open list ordinata + Closed list & Solo stack di ricorsione \\
Cicli & Gestiti dalla Closed list & Gestiti con path checking \\
Ottimale & Sì (con $h$ ammissibile) & Sì (con $h$ ammissibile) \\
Completo & Sì & Sì (con costi $> 0$) \\
Statistiche & Nodi espansi + max open & Numero iterazioni \\
\bottomrule
\end{tabular}
\end{center}

% ======================================================================
% PARTE II - ASP/CLINGO
% ======================================================================
\part{Answer Set Programming (Clingo)}

% =====================================================================
\section{Fondamenti di ASP}
% =====================================================================

\subsection{Paradigma}

L'ASP è un paradigma \textbf{dichiarativo}: si descrivono i \textbf{vincoli} che
la soluzione deve soddisfare. Il solver (Clingo) trova gli \textbf{answer set}
(modelli stabili) che soddisfano tutti i vincoli.

\textbf{Pipeline}:
\begin{center}
\texttt{Grounding} (variabili $\to$ costanti) $\longrightarrow$ \texttt{Solving} (cerca modelli stabili) $\longrightarrow$ \texttt{Ottimizzazione}
\end{center}

\subsection{Costrutti Fondamentali}

\begin{center}
\renewcommand{\arraystretch}{1.4}
\begin{tabular}{p{3.2cm}p{4.2cm}p{6cm}}
\toprule
\textbf{Costrutto} & \textbf{Sintassi} & \textbf{Significato} \\
\midrule
Fatto & \texttt{p(a).} & \texttt{p(a)} è sempre vero \\
Regola normale & \texttt{h :- b1, b2.} & derivazione condizionale \\
Negazione NAF & \texttt{not p} & vero se \texttt{p} non è nell'answer set \\
\textbf{Choice rule} & \texttt{\{h\} :- body.} & \texttt{h} può o non può essere incluso \\
\textbf{Vincolo hard} & \texttt{:- body.} & elimina modelli dove body è vero \\
Aggregato count & \texttt{\#count\{x : p(x)\}} & conta gli x con p(x) \\
Aggregato sum & \texttt{\#sum\{w,k : p(k,w)\}} & somma dei pesi \\
Modulo & \texttt{N \textbackslash M} & resto intero di N/M \\
Minimizzazione & \texttt{\#minimize\{w,k : p\}.} & minimizza la somma \\
Show & \texttt{\#show p/n.} & mostra solo \texttt{p} nell'output \\
\bottomrule
\end{tabular}
\end{center}

\begin{defbox}{Choice Rule: la distinzione chiave da Prolog}
In Prolog, se il corpo è vero la testa \emph{deve} essere vera.
In ASP, \texttt{\{h\} :- body} significa che \texttt{h} \emph{può} essere
vero se body è vero, ma non è obbligatorio. Il solver esplora entrambe le
possibilità e mantiene i modelli che soddisfano i vincoli.
\end{defbox}

\begin{defbox}{Vincolo Hard: il filtro}
\texttt{:- body} elimina dall'insieme di answer set tutti i modelli dove
\texttt{body} è vero. Equivale a dire: ``è \textbf{vietato} che body sia
soddisfatto''.
\end{defbox}

% =====================================================================
\section{Progetto Torneo Pokémon -- \texttt{torneo.lp}}
% =====================================================================

\subsection{Sistema Svizzero}

Regole di abbinamento:
\begin{enumerate}[noitemsep]
  \item Giocatori con punteggio simile si affrontano (spareggio).
  \item Nessun rematching (stessa coppia non si affronta due volte).
  \item Se numero dispari: un giocatore riceve il BYE (vittoria auto). Nessuno lo prende due volte.
\end{enumerate}

\subsection{Normalizzazione: \texttt{pid} e \texttt{played\_already}}

\begin{lstlisting}[language=ASP]
% Simmetria degli incontri
played_already(P1, P2) :- played(P1, P2).
played_already(P1, P2) :- played(P2, P1).

% Indice canonico: quanti giocatori vengono prima di P lessicograficamente
pid(P, I) :- player(P),
             I = #count{ Q : player(Q), Q != P, Q < P }.
\end{lstlisting}

\textbf{Scopo di \texttt{pid}}: assegnare un ordine totale ai giocatori per
generare ogni coppia \{P1, P2\} in una sola forma canonica (solo quando
$\texttt{pid}(P1) < \texttt{pid}(P2)$), evitando di avere sia
\texttt{match(A,B)} che \texttt{match(B,A)}.

\subsection{Generazione}

\begin{lstlisting}[language=ASP]
% match solo in forma canonica (pid(P1) < pid(P2))
{ match(P1, P2) } :-
    player(P1), player(P2),
    pid(P1, I1), pid(P2, I2), I1 < I2,
    not played_already(P1, P2).

% BYE solo per chi non lo ha mai avuto
{ bye(P) } :- player(P), not had_bye(P).
\end{lstlisting}

\subsection{Vincoli Hard}

\begin{lstlisting}[language=ASP]
in_match(P) :- match(P, _).
in_match(P) :- match(_, P).

% Ogni giocatore: esattamente match OR bye
:- player(P), not in_match(P), not bye(P).  % deve fare qualcosa
:- player(P), in_match(P), bye(P).           % non entrambi

% Al massimo in un match (3 casi di conflitto)
:- match(P, Q1), match(P, Q2), Q1 != Q2.    % P in due match come P1
:- match(P1, Q), match(P2, Q), P1 != P2.    % Q in due match come P2
:- match(P, _),  match(_, P).               % P in entrambi i ruoli

% Un solo BYE per turno
:- 2 <= #count{ P : bye(P) }.

% Parita': BYE iff numero dispari
player_count(N) :- N = #count{ P : player(P) }.
:- player_count(N), 0 == N \ 2, bye(_).     % pari -> no BYE
:- player_count(N), 1 == N \ 2, not bye(_). % dispari -> BYE obbligatorio
\end{lstlisting}

\subsection{Ottimizzazione}

\begin{lstlisting}[language=ASP]
score_diff(P1, P2, D) :-
    match(P1, P2), score(P1, S1), score(P2, S2),
    S1 >= S2, D = S1 - S2.
score_diff(P1, P2, D) :-
    match(P1, P2), score(P1, S1), score(P2, S2),
    S2 > S1, D = S2 - S1.

#minimize { D, P1, P2 : score_diff(P1, P2, D) }.
\end{lstlisting}

Minimizza $\sum |S_{P1} - S_{P2}|$ su tutti i match. Match tra giocatori con
punteggio uguale contribuiscono 0 alla somma (preferiti).

\subsection{Istanze}

\begin{center}
\begin{tabular}{lllll}
\toprule
\textbf{Turno} & \textbf{Giocatori} & \textbf{Parità} & \textbf{BYE} & \textbf{Ottimiz.} \\
\midrule
Turno 1 & 5 (dispari) & 5 mod 2 = 1 & Obbligatorio & Tutti a 0: qualsiasi match \\
Turno 2 & 8 (pari) & 8 mod 2 = 0 & Vietato & Abbina i "2" tra loro \\
\bottomrule
\end{tabular}
\end{center}

% ======================================================================
% PARTE III - CLIPS
% ======================================================================
\part{CLIPS -- Expert System}

% =====================================================================
\section{Fondamenti di CLIPS}
% =====================================================================

\subsection{Cos'è CLIPS}

CLIPS (C Language Integrated Production System) è un framework per
\textbf{sistemi esperti} basato su regole di produzione.

\begin{center}
\begin{tabular}{lll}
\toprule
\textbf{Componente} & \textbf{Nome CLIPS} & \textbf{Funzione} \\
\midrule
Fatti & Working Memory (WM) & Ciò che il sistema ``sa'' \\
Regole & Production Rules & \texttt{if} condizioni \texttt{then} azioni \\
Motore & Inference Engine & Seleziona e spara le regole \\
\bottomrule
\end{tabular}
\end{center}

\subsection{Ciclo di Inferenza}

\begin{enumerate}[noitemsep]
  \item \textbf{Match}: trova tutte le regole le cui condizioni (LHS) sono soddisfatte dai fatti in WM $\to$ \emph{agenda}.
  \item \textbf{Conflict resolution}: seleziona la regola con \textbf{salience} maggiore.
  \item \textbf{Fire}: esegue le azioni (RHS) della regola (assert, retract, modify, focus, ...).
  \item Ripete finché l'agenda è vuota.
\end{enumerate}

\subsection{Costrutti Base}

\begin{lstlisting}[language=CLIPS]
;;; Template (struttura dati)
(deftemplate MAIN::cell
    (slot x       (type INTEGER))
    (slot visited (type SYMBOL) (default no))
    (slot pit-cf  (type FLOAT)  (default 0.0))
)

;;; Regola
(defrule MAIN::nome-regola
    (declare (salience 10))    ; priorita' (default 0)
    (percept (x ?x) (y ?y))   ; condizioni (LHS)
    ?c <- (cell (x ?x))       ; bind del fatto a ?c
=>
    (modify ?c (visited yes))  ; azione: modifica ?c
    (assert (nuova-cosa))      ; aggiungi un fatto
)

;;; Modulo
(defmodule PERCEPT (import MAIN ?ALL) (export ?ALL))
\end{lstlisting}

\begin{defbox}{Salience}
La salience determina l'ordine di esecuzione delle regole. Range: $-10000$ a
$+10000$, default 0. Regole con salience più alta sparano prima.
\end{defbox}

\begin{defbox}{Moduli e Focus}
Un modulo è un namespace per regole. \texttt{(focus MODULO)} mette il modulo
in cima allo stack: solo le sue regole possono sparare. Quando il modulo esaurisce
le regole attive, si torna al modulo precedente nello stack.
\end{defbox}

% =====================================================================
\section{Progetto Wumpus -- \texttt{agent.clp} + \texttt{wumpus.py}}
% =====================================================================

\subsection{Architettura Generale}

\textbf{Python} simula il mondo (griglia, Wumpus, pozzi, oro) e fa da bridge:
\begin{enumerate}[noitemsep]
  \item Calcola le percezioni dalla posizione corrente.
  \item Inietta un fatto \texttt{(percept ...)} in CLIPS.
  \item Chiama \texttt{env.run()} per far sparare le regole.
  \item Legge il fatto \texttt{(action ...)} e applica l'azione al mondo.
\end{enumerate}

\subsection{Moduli dell'Agente CF}

\begin{center}
\begin{tikzpicture}[
  box/.style={draw,rounded corners,minimum width=2.5cm,minimum height=0.9cm,align=center,font=\small},
  arrow/.style={-{Stealth},thick}
]
  \node[box,fill=gray!20]   (main)   {MAIN\\{\tiny Coordinamento}};
  \node[box,fill=blue!20, right=1.8cm of main]   (perc)   {PERCEPT\\{\tiny Aggiorna stato}};
  \node[box,fill=orange!20, right=1.8cm of perc]  (infer)  {INFER\\{\tiny Certainty Factors}};
  \node[box,fill=green!20, right=1.8cm of infer]  (decide) {DECIDE\\{\tiny Sceglie azione}};
  \draw[arrow] (main) -- node[above,font=\tiny]{\texttt{focus}} (perc);
  \draw[arrow] (perc) -- node[above,font=\tiny]{\texttt{focus}} (infer);
  \draw[arrow] (infer) -- node[above,font=\tiny]{\texttt{focus}} (decide);
\end{tikzpicture}
\end{center}

\subsection{Template Principali}

\begin{lstlisting}[language=CLIPS]
(deftemplate MAIN::cell
    (slot x) (slot y)
    (slot visited   (type SYMBOL) (default no))
    (slot safe      (type SYMBOL) (default unknown))  ; yes/no/unknown
    (slot pit-cf    (type FLOAT)  (default 0.0))      ; in [-1, +1]
    (slot wumpus-cf (type FLOAT)  (default 0.0))
)
(deftemplate MAIN::percept   ; iniettato da Python ogni ciclo
    (slot x) (slot y) (slot steps) (slot energy)
    (slot breeze (type SYMBOL))   ; yes/no
    (slot stench (type SYMBOL))   ; yes/no
    (slot glitter (type SYMBOL))  ; yes/no
)
(deftemplate MAIN::pending-action
    (slot move (type SYMBOL))
    (slot priority (type INTEGER) (default 0))
)
(deftemplate MAIN::cf-done  ; impedisce loop da modify
    (slot x) (slot y) (slot kind)  ; kind = pit o wumpus
)
\end{lstlisting}

\subsection{Certainty Factors (CF)}

\begin{mathbox}{Teoria dei Certainty Factors}
Un CF $\in [-1, +1]$ quantifica la certezza di un'ipotesi:
\begin{itemize}[noitemsep]
  \item $-1$: certezza assoluta che sia \textbf{falsa}
  \item $0$: completa ignoranza
  \item $+1$: certezza assoluta che sia \textbf{vera}
\end{itemize}

\textbf{Evidenza positiva} (brezza/odore rilevato, $\alpha = 0.5$):
\[
  CF_{\text{new}} = CF + (1 - CF) \times 0.5
\]

\textbf{Evidenza negativa} (nessuna brezza/odore, $\alpha = 0.7$):
\[
  CF_{\text{new}} = CF + (-1 - CF) \times 0.7
\]

Le formule hanno rendimento decrescente: il CF non supera mai $\pm 1$.
Partendo da 0, dopo 1 evidenza positiva: $0.5$; dopo 2: $0.75$; dopo 3: $0.875$.
Dopo 1 evidenza negativa: $-0.7$; dopo 2: $-0.91$.
\end{mathbox}

\subsubsection{Esempio di regola CF -- INFER}

\begin{lstlisting}[language=CLIPS]
(defrule INFER::breeze-north
    (declare (salience 10))
    (percept (x ?x) (y ?y) (breeze yes))
    ?c <- (cell (x ?cx) (y ?cy) (visited no) (pit-cf ?cf))
    (test (= ?cx ?x))
    (test (= ?cy (+ ?y 1)))                     ; cella a nord
    (not (cf-done (x ?cx) (y ?cy) (kind pit)))  ; non gia' aggiornata
=>
    (assert (cf-done (x ?cx) (y ?cy) (kind pit)))         ; blocca ri-attivazione
    (modify ?c (pit-cf (+ ?cf (* (- 1.0 ?cf) 0.5))))     ; CF_new = CF + (1-CF)*0.5
)
\end{lstlisting}

\begin{warnbox}{Il problema del loop da \texttt{modify}}
\texttt{(modify ?f ...)} in CLIPS fa retract + re-assert del fatto, rendendo il
fatto ``nuovo'' e riattivando tutte le regole che lo matchano.
\textbf{Senza} \texttt{cf-done}: la regola sparerebbe infinite volte sulla stessa
cella. \textbf{Con} \texttt{cf-done}: la condizione \texttt{(not (cf-done ...))}
blocca la ri-attivazione nello stesso ciclo.
\end{warnbox}

\subsubsection{Classificazione delle Celle}

\begin{lstlisting}[language=CLIPS]
(defrule INFER::mark-safe     ; pit-cf < -0.5 E wumpus-cf < -0.5 -> sicura
    ?c <- (cell (visited no) (safe unknown) (pit-cf ?p) (wumpus-cf ?w))
    (test (< ?p -0.5)) (test (< ?w -0.5))
=> (modify ?c (safe yes)) )

(defrule INFER::mark-dangerous-pit   ; pit-cf > 0.6 -> pericolosa
    ?c <- (cell (visited no) (safe unknown) (pit-cf ?p))
    (test (> ?p 0.6))
=> (modify ?c (safe no)) )
\end{lstlisting}

\subsection{Gerarchia delle Azioni -- DECIDE}

\begin{center}
\renewcommand{\arraystretch}{1.3}
\begin{tabular}{cll}
\toprule
\textbf{Priorità} & \textbf{Azione} & \textbf{Quando} \\
\midrule
100 & \texttt{grab} & glitter presente \\
90  & \texttt{quit} & energia $\le 30$ \\
80  & \texttt{up/down/left/right} & cella \texttt{safe=yes} non visitata adiacente \\
60  & \texttt{pit-north/...} & CF pit $> 0.3$ e CF wumpus $< 0.3$ (indovina) \\
10  & \texttt{up/...} (esplora) & CF pit $< 0.5$ e CF wumpus $< 0.5$ \\
5   & \texttt{up/...} (coraggioso) & cella non esplicitamente pericolosa \\
0   & \texttt{quit} (fallback) & nessuna altra opzione \\
\bottomrule
\end{tabular}
\end{center}

\textbf{Meccanismo di selezione}: le regole asseriscono \texttt{pending-action}
con una priorità. Poi \texttt{cleanup-lower} rimuove quelle subottimali e
\texttt{emit-action} promuove la vincente ad \texttt{action}.

\begin{lstlisting}[language=CLIPS]
(defrule DECIDE::cleanup-lower
    (declare (salience 0))
    ?pa <- (pending-action (priority ?p1))
    (pending-action (priority ?p2))
    (test (> ?p2 ?p1))
=> (retract ?pa) )

(defrule DECIDE::emit-action
    (declare (salience -1))
    (not (action))
    ?pa <- (pending-action (move ?m) (priority ?p))
    (not (pending-action (priority ?p2&:(> ?p2 ?p))))
=> (assert (action (move ?m))) (retract ?pa) )
\end{lstlisting}

\subsection{Agente Greedy vs Agente CF}

\begin{center}
\renewcommand{\arraystretch}{1.4}
\begin{tabular}{lll}
\toprule
\textbf{Aspetto} & \textbf{agent.clp (CF)} & \textbf{agent\_greedy.clp} \\
\midrule
Moduli & 4 (MAIN/PERCEPT/INFER/DECIDE) & 1 (MAIN) \\
Cella pericolosa & CF $\in [-1, +1]$, graduabile & \texttt{dangerous = yes/no} (binario) \\
Aggiornamento & Incrementale, reversibile & Irreversibile (1 brezza = pericolo) \\
Pit-guess & Sì (se CF moderato) & No \\
Blocco & Raramente (ha fallback multipli) & Spesso (tutte le celle \texttt{dangerous}) \\
\bottomrule
\end{tabular}
\end{center}

% ======================================================================
% PARTE IV - SCHEMA RIASSUNTIVO
% ======================================================================
\part{Schema Riassuntivo per l'Esame}

% =====================================================================
\section{Confronto tra i Paradigmi}
% =====================================================================

\begin{center}
\renewcommand{\arraystretch}{1.4}
\begin{tabular}{p{3cm}p{3.8cm}p{3.8cm}p{3.8cm}}
\toprule
\textbf{Aspetto} & \textbf{Prolog} & \textbf{ASP (Clingo)} & \textbf{CLIPS} \\
\midrule
Paradigma & Logico / dichiarativo & Logico / constraint & Regole di produzione \\
Negazione & NAF (\texttt{\textbackslash+}) & NAF (\texttt{not}) & Patterns negativi (\texttt{not}) \\
Scelta & Backtracking & Choice rules \texttt{\{h\}} & Agenda (salience) \\
Vincoli & Clausole Horn & Integrity constraints & Condizioni LHS \\
Ottimizzazione & No (built-in) & \texttt{\#minimize/maximize} & No (built-in) \\
Uso principale & Ricerca, logica, NLP & Combinatoria, pianificazione & Sistemi esperti \\
\bottomrule
\end{tabular}
\end{center}

% =====================================================================
\section{Domande di Esame Tipiche}
% =====================================================================

\subsection{Su Prolog}

\begin{enumerate}
  \item \textbf{Cos'è il backtracking e quando si attiva?}
        Quando un goal fallisce o si vuole la prossima soluzione, Prolog torna
        all'ultimo punto di scelta e prova l'alternativa successiva.

  \item \textbf{Differenza tra \texttt{=} e \texttt{is}?}
        \texttt{X = 3+2} unifica X con il termine \texttt{+(3,2)} (non valuta).
        \texttt{X is 3+2} valuta l'espressione e unifica X con 5.

  \item \textbf{Cosa fa il cut (\texttt{!})?}
        Elimina i punti di backtracking: impedisce di tornare indietro sulla
        clausola corrente e sulle sue alternative.

  \item \textbf{Differenza tra cut verde e rosso?}
        Verde: solo efficienza, non cambia il significato. Rosso: cambia il
        significato, elimina soluzioni che senza cut sarebbero trovate.

  \item \textbf{Perché \texttt{findall} non fallisce mai?}
        Perché se il goal non ha soluzioni restituisce \texttt{[]}.

  \item \textbf{A* è ottimale? Quando?}
        Sì, se l'euristica è \textbf{ammissibile}: $h(n) \le h^*(n)$ per ogni $n$.

  \item \textbf{Vantaggio principale di IDA* su A*?}
        Memoria lineare $O(b \cdot d)$ invece di esponenziale $O(b^d)$.

  \item \textbf{Perché la distanza Manhattan è ammissibile?}
        Conta i passi minimi senza ostacoli. I muri possono solo aumentare il
        costo reale, quindi non sovrastima mai.
\end{enumerate}

\subsection{Su ASP/Clingo}

\begin{enumerate}
  \item \textbf{Differenza tra regola normale e choice rule?}
        Normale (\texttt{h :- body}): se body è vero, h \textbf{deve} essere vero.
        Choice (\texttt{\{h\} :- body}): se body è vero, h \textbf{può} essere vero.

  \item \textbf{Cosa fa un integrity constraint?}
        \texttt{:- body} elimina tutti i modelli dove body è vero.

  \item \textbf{Perché si usa \texttt{pid} nel torneo?}
        Per generare ogni coppia di giocatori in una sola forma canonica
        (evitare sia \texttt{match(A,B)} sia \texttt{match(B,A)}).

  \item \textbf{Cosa minimizza \texttt{\#minimize}?}
        La somma dei pesi degli elementi dell'insieme che soddisfano la condizione.
        Nel torneo: somma delle differenze di punteggio tra gli abbinati.

  \item \textbf{Come si gestisce la parità nel torneo?}
        Con l'operatore modulo \texttt{N \textbackslash 2}: se resto 0 (pari) si
        vieta il BYE; se resto 1 (dispari) si impone il BYE.
\end{enumerate}

\subsection{Su CLIPS}

\begin{enumerate}
  \item \textbf{Cos'è la salience?}
        Un valore intero che determina la priorità di una regola. Valori più alti
        sparano prima. Default: 0.

  \item \textbf{Perché si usa \texttt{cf-done}?}
        Per evitare loop infiniti causati da \texttt{modify}: ogni \texttt{modify}
        ri-asserta il fatto rendendolo nuovo, riattivando le regole. Il fatto
        \texttt{cf-done} blocca la ri-attivazione nello stesso ciclo.

  \item \textbf{Cosa sono i Certainty Factors?}
        Valori in $[-1, +1]$ che quantificano la certezza di un'ipotesi.
        Si aggiornano con formule che combinano evidenza positiva e negativa.

  \item \textbf{Perché l'agente CF è superiore al greedy?}
        Perché gestisce l'incertezza in modo graduato: può rivalutare una cella
        parzialmente pericolosa se nuove evidenze negative la abbassano.
        Il greedy è irreversibile: una singola brezza marca la cella come
        pericolosa per sempre.

  \item \textbf{Perché l'action si rimuove prima di iniettare il percept?}
        Perché la regola \texttt{handle-percept} si attiva solo quando
        \texttt{(not (action))}. Se l'action del ciclo precedente rimane
        in WM, il nuovo ciclo non parte.
\end{enumerate}

% =====================================================================
\section{Formule da Ricordare}
% =====================================================================

\begin{center}
\renewcommand{\arraystretch}{1.6}
\begin{tabular}{lll}
\toprule
\textbf{Formula} & \textbf{Contesto} & \textbf{Significato} \\
\midrule
$f(n) = g(n) + h(n)$ & A* / IDA* & costo totale stimato attraverso $n$ \\
$h(n) \le h^*(n)$ & Ammissibilità & condizione per ottimalità \\
$h_{\text{Manhattan}} = |r_1-r_2| + |c_1-c_2|$ & Labirinto / Puzzle & stima del percorso \\
$CF_{\text{new}} = CF + (1-CF) \times 0.5$ & CLIPS / CF positivo & aggiornamento con brezza \\
$CF_{\text{new}} = CF + (-1-CF) \times 0.7$ & CLIPS / CF negativo & aggiornamento senza brezza \\
$N \bmod 2 = 0 \Rightarrow$ no BYE & ASP Torneo & vincolo di parità \\
\bottomrule
\end{tabular}
\end{center}

\end{document}
